\documentclass[11pt]{article}

\usepackage[margin=1in]{geometry}
\usepackage{amsmath,amssymb}
\usepackage{algorithm}
\usepackage{algpseudocode}
\usepackage{hyperref}

\title{Group-Delay Equalization for Complex IIR Filters:\\ Design Procedure}
\author{}
\date{}

\begin{document}
\maketitle

\noindent This section states the complete equalizer design procedure on
its own, in a form intended to be reusable close to verbatim in a paper.

\section{Filter model and all-pass building block}

Let $H(z)$ be a (possibly complex, non-conjugate-symmetric) discrete-time
IIR filter with passband $[f_{p1}, f_{p2}]$ in normalized cyclic
frequency ($F_s = 1$), and let $\tau_H(f)$ denote its group delay.
Complex IIR filters designed without an arithmetic-symmetry constraint
typically exhibit a markedly non-flat, ``bathtub''-shaped group delay:
two peaks near the passband edges bridged by a lower floor across the
interior. The goal is a cascaded all-pass equalizer that flattens
$\tau_H$ across the passband without altering $|H(e^{j\omega})|$.

The equalizer is a cascade of $N$ first-order discrete all-pass
sections. Each section is parameterized by a single complex pole
$p_i = r_i e^{j\theta_i}$, $0 < r_i < 1$:
\begin{equation}
A_i(z) \;=\; \frac{z^{-1} - \overline{p_i}}{1 - p_i\, z^{-1}}.
\label{eq:allpass}
\end{equation}
This construction is magnitude-preserving for any complex $p_i$:
$|A_i(e^{j\omega})| = 1$ for all real $\omega$, so cascading $A_i$ onto
$H$ leaves $|H(e^{j\omega})|$ unchanged and only modifies phase/group
delay. The group-delay contribution of a single section is the Poisson
kernel
\begin{equation}
\tau_i(\omega) \;=\; \frac{1 - r_i^2}{1 - 2 r_i \cos(\omega - \theta_i) + r_i^2},
\qquad \tau_i(\omega) > 0 \ \ \forall\, \omega.
\label{eq:poisson}
\end{equation}
Because $\tau_i(\omega) > 0$ everywhere, an all-pass cascade can only
\emph{add} delay: it can raise a low floor toward the level of existing
peaks, but it cannot directly lower a peak, and a section placed too
close to an existing peak only adds to it.

For implementation convenience, sections are grouped into $M$
\emph{clusters}: cluster $i$ consists of $c_i$ identical stacked
sections sharing one pole $(r_i, \theta_i)$, so $N = \sum_i c_i$, and the
total equalizer delay is
\begin{equation}
\tau_{eq}(\omega; \boldsymbol\theta, \mathbf r) \;=\; \sum_{i=1}^{M} c_i\, \tau_i(\omega; r_i, \theta_i),
\qquad
\tau(\omega) \;=\; \tau_H(\omega) + \tau_{eq}(\omega).
\label{eq:total}
\end{equation}

\section{Design objective: anchored, gain-weighted equi-ripple balance}

Fix an \textbf{anchor} level once, from the \emph{unequalized} filter
alone:
\begin{equation}
D_0 \;=\; \operatorname{mean}\Big(\tau_H(f) : f \text{ a local extremum of } \tau_H\Big),
\label{eq:anchor}
\end{equation}
in practice the mean of the filter's own two edge-of-passband peaks.
$D_0$ is never recomputed once fixed, which keeps the objective anchored
to an external, fixed target rather than to the equalizer's own
(evolving) output -- an unanchored objective was found in practice to
let the optimizer invent new peaks and match them to each other at an
arbitrarily poor common level.

For a given $(\boldsymbol\theta, \mathbf r)$, let $\{f_k\}_{k=1}^K$ be
the frequencies of \emph{every} local extremum (maxima and minima) of
$\tau(\cdot)$ over the passband expanded by a margin fraction $\mu$ on
each side (default $\mu=0.10$), redetected from scratch at every
iteration (Section~\ref{sec:iteration}). Define the gain-weighted
deviation from the anchor at each extremum:
\begin{equation}
w_k \;=\; \big|\tau(f_k) - D_0\big| \cdot |H(f_k)|, \qquad k = 1, \dots, K.
\label{eq:wk}
\end{equation}
The design goal is the classical Chebyshev equi-ripple condition: choose
$(\boldsymbol\theta, \mathbf r)$ so that all $w_k$ are equal. Weighting
by the filter's own gain $|H(f_k)|$ allows an extremum sitting where
$|H|$ has already rolled off to carry a proportionally larger raw delay
deviation without violating equality of the weighted terms -- the
deviation that matters is the one seen at the filter's output, not the
raw group-delay number.

\section{One Newton/least-squares iteration}
\label{sec:iteration}

Given a current configuration $(\boldsymbol\theta, \mathbf r)$, one
iteration proceeds in two steps.

\paragraph{Step 1 -- extremum detection and refinement.} Evaluate
$\tau''(\omega)$ on a uniform grid over the expanded band and locate its
zero crossings (each pair of consecutive crossings, plus the two
boundary segments, brackets at most one true extremum of $\tau$, since
$\tau'$ is monotonic wherever $\tau''$ keeps one sign). Within each
occupied bracket, confirm the extremum type from the sign of $\tau'$ at
the bracket's endpoints, then refine its location with one Newton step:
\begin{equation}
f_k \;\leftarrow\; f_k \;-\; \frac{\tau'(f_k)}{\tau''(f_k)}.
\label{eq:newtonrefine}
\end{equation}

\paragraph{Step 2 -- parameter update.} Because cluster $i$'s
contribution depends only on its own $(r_i,\theta_i)$, the Jacobian of
$\tau$ with respect to each cluster parameter is available in closed
form:
\begin{align}
\frac{\partial \tau(\omega)}{\partial \theta_i} &= c_i\,\frac{2 r_i (1-r_i^2)\sin(\omega-\theta_i)}{D_i(\omega)^2}, \label{eq:jactheta}\\[4pt]
\frac{\partial \tau(\omega)}{\partial r_i} &= c_i\left(\frac{-2 r_i}{D_i(\omega)} \;-\; \frac{2(1-r_i^2)\big(r_i - \cos(\omega-\theta_i)\big)}{D_i(\omega)^2}\right), \label{eq:jacr}\\[4pt]
D_i(\omega) &= 1 - 2 r_i \cos(\omega-\theta_i) + r_i^2. \label{eq:Di}
\end{align}
Assemble the Jacobian $J \in \mathbb{R}^{K \times M}$ (angle-only) or
$J \in \mathbb{R}^{K \times 2M}$ (joint angle+radius) of the weighted
deviations, row $k$ given by
$J_{k,i} = \operatorname{sign}\big(\tau(f_k)-D_0\big)\,|H(f_k)|\,\partial\tau(f_k)/\partial\theta_i$
(and analogously for $\partial/\partial r_i$). Solve, via the
Moore--Penrose pseudoinverse (the system is neither reliably square nor
consistently over/under-determined, since $K$ changes call to call),
\begin{equation}
\Delta \;=\; J^{+}\big(\bar w \,\mathbf 1 - \mathbf w\big), \qquad \bar w = \frac1K \sum_{k=1}^K w_k,
\label{eq:lstsq}
\end{equation}
the least-squares parameter step that drives every $w_k$ toward their
common mean, then apply a damped update
\begin{equation}
\boldsymbol\theta \leftarrow \boldsymbol\theta + \alpha\, \Delta_\theta, \qquad
\mathbf r \leftarrow \mathbf r + \alpha\, \Delta_r \ \ (\text{if radii are free}),
\label{eq:damped}
\end{equation}
with damping factor $\alpha \in (0,1]$, and radii clamped to
$[r_{\min}, r_{\max}] = [0.3, 0.995]$ after the update.

\section{Two-phase iteration schedule}

\textbf{Phase 1 (angle-only).} Radii held fixed; repeat
Section~\ref{sec:iteration} with damping $\alpha = \alpha_1$ until the
ripple spread $S = \max_k w_k - \min_k w_k$ changes by less than a
tolerance $\epsilon$ between consecutive iterations, or a maximum
iteration budget is reached. This phase converges monotonically in
practice.

\textbf{Phase 2 (joint angle + radius).} Continuing from the Phase~1
result, free the radii and repeat Section~\ref{sec:iteration} with
damping $\alpha = \alpha_2$, now updating $J$'s $r$-columns as well.
This phase does \emph{not} converge monotonically: $S$ reaches a minimum
after a small number of steps, then drifts upward again while flatness
just outside the nominal passband continues to improve -- a genuine
trade-off, not simple convergence to a fixed point. Accordingly, Phase~2
tracks the best (smallest) $S$ seen so far and stops after $P$
consecutive non-improving iterations, \textbf{returning the best
configuration found}, not the final iterate.

\section{Starting configuration and step-size selection}

The iteration is initialized with $M$ clusters of $c$ stacked sections
each, pole angles spread across the passband (evenly, or symmetrically
if the filter's own specification is symmetric), and a common starting
radius $r_0$. Two empirical rules govern the free parameters:

\begin{itemize}
\item \textbf{Damping must shrink as $M$ grows relative to $K$.} With
  more free angle/radius parameters against roughly the same number of
  tracked extrema, the pseudoinverse step becomes more aggressive per
  unit residual; the default $\alpha=0.5$, stable at $M=3$, produces
  outright divergence at $M=5$ (oscillating extremum count, radii
  pinned at their clamp bounds) unless $\alpha_1,\alpha_2$ are reduced
  (e.g.\ to $0.2/0.1$, or smaller still when perturbing a starting point
  that is already known to be reasonable).
\item \textbf{Lower $r_0$ suits wider passbands.} A single section's own
  peak height $(1+r_i)/(1-r_i)$ grows far faster with $r_i$ than its
  fractional leakage to distant frequencies shrinks; for a passband
  spanning a large fraction of the full cycle, a lower, broader starting
  radius blends clusters into a smooth correction rather than a
  sequence of narrow, high-$Q$ bumps.
\end{itemize}

\section{Algorithm summary}

\begin{algorithm}[h]
\caption{Anchored equi-ripple group-delay equalization}
\begin{algorithmic}[1]
\Require $H(z)$; passband $[f_{p1},f_{p2}]$; expansion margin $\mu$;
  cluster count $M$; stages per cluster $c_i$; starting angles
  $\boldsymbol\theta^{(0)}$ and radius $r_0$; damping $\alpha_1,\alpha_2$;
  tolerance $\epsilon$; patience $P$.
\State Compute $\tau_H$ over the expanded band; set anchor $D_0$ from
  $\tau_H$'s own extrema (Eq.~\ref{eq:anchor}).
\State $\boldsymbol\theta \gets \boldsymbol\theta^{(0)}$; \quad
  $\mathbf r \gets r_0\mathbf 1$.
\Statex \textbf{Phase 1 (angle-only):}
\Repeat
  \State Update $\boldsymbol\theta$ via Section~\ref{sec:iteration} with
    radii fixed, damping $\alpha_1$.
\Until{$|S^{(t)}-S^{(t-1)}| < \epsilon$}
\Statex \textbf{Phase 2 (joint angle+radius):}
\State $S_{\text{best}} \gets S$; \quad $(\boldsymbol\theta_{\text{best}}, \mathbf r_{\text{best}}) \gets (\boldsymbol\theta, \mathbf r)$; \quad $\text{noImprove} \gets 0$
\Repeat
  \State Update $(\boldsymbol\theta,\mathbf r)$ via
    Section~\ref{sec:iteration} with radii free, damping $\alpha_2$.
  \If{$S < S_{\text{best}}$}
    \State $S_{\text{best}} \gets S$; \quad $(\boldsymbol\theta_{\text{best}}, \mathbf r_{\text{best}}) \gets (\boldsymbol\theta, \mathbf r)$; \quad $\text{noImprove} \gets 0$
  \Else
    \State $\text{noImprove} \gets \text{noImprove} + 1$
  \EndIf
\Until{$\text{noImprove} \geq P$}
\State \Return equalizer $\{(r_i,\theta_i,c_i)\}_{i=1}^M = (\boldsymbol\theta_{\text{best}}, \mathbf r_{\text{best}}, \mathbf c)$
\end{algorithmic}
\end{algorithm}

\end{document}
